\documentclass[11pt]{article}
\usepackage{amsmath,amssymb,amsthm}
\usepackage{algorithm,algorithmic}
\usepackage{hyperref}
\usepackage{xcolor}
\newtheorem{theorem}{Theorem}
\newtheorem{lemma}{Lemma}
\newtheorem{assumption}{Assumption}
\newcommand{\norm}[1]{\left\lVert#1\right\rVert}
\newcommand{\ip}[2]{\left\langle #1,#2\right\rangle}
\newcommand{\grad}{\nabla}
\begin{document}

%=====================================================================
\section*{Heavy‑Ball Momentum (HB)}
%=====================================================================

\section*{Mathematical Formulation}
Let $f:\mathbb R^{d}\rightarrow\mathbb R$ be differentiable.  
Given step size $\alpha>0$ and momentum parameter $\beta\in[0,1)$, the
Heavy‑Ball iteration is  

\[
\boxed{
\begin{aligned}
x_{k+1}&=x_{k}+\beta\bigl(x_{k}-x_{k-1}\bigr)-\alpha\,\grad f(x_{k}),\qquad k\ge 0,
\end{aligned}}
\tag{HB}
\]

with an initial pair $(x_{0},x_{-1})$ (commonly $x_{-1}=x_{0}$).

\section*{Pseudocode}
\begin{algorithm}[H]
\caption{Heavy‑Ball Momentum}
\begin{algorithmic}[1]
\REQUIRE step size $\alpha>0$, momentum $\beta\in[0,1)$,
          initial point $x_{0}\in\mathbb R^{d}$,
          (optional) $x_{-1}=x_{0}$.
\FOR{$k=0,1,2,\dots$}
    \STATE $g_{k}\gets\grad f(x_{k})$ \COMMENT{gradient at current iterate}
    \STATE $x_{k+1}\gets x_{k}+\beta\bigl(x_{k}-x_{k-1}\bigr)-\alpha\,g_{k}$ \label{line:update}
    \IF{stopping criterion satisfied}
        \STATE \textbf{break}
    \ENDIF
\ENDFOR
\RETURN $x_{k+1}$
\end{algorithmic}
\end{algorithm}

\section*{Dry Run Example}
Consider the one–dimensional quadratic  

\[
f(w)=(w-3)^{2},\qquad\grad f(w)=2(w-3).
\]

Choose $\alpha=0.1$, $\beta=0.5$, and initialise $w_{0}=0$, $w_{-1}=0$.

\begin{center}
\begin{tabular}{c|c|c|c}
$k$ & $g_{k}=2(w_{k}-3)$ & $w_{k+1}=w_{k}+\beta(w_{k}-w_{k-1})-\alpha g_{k}$ & $f(w_{k})$\\\hline
0 & $2(0-3)=-6$ & $0+0.5(0-0)-0.1(-6)=0.6$ & $(0.6-3)^{2}=5.76$\\
1 & $2(0.6-3)=-4.8$ & $0.6+0.5(0.6-0)-0.1(-4.8)=1.08$ & $(1.08-3)^{2}=3.6864$\\
2 & $2(1.08-3)=-3.84$ & $1.08+0.5(1.08-0.6)-0.1(-3.84)=1.452$ & $(1.452-3)^{2}=2.3990$
\end{tabular}
\end{center}
The iterates rapidly approach the minimiser $w^{*}=3$.

%=====================================================================
\section*{Assumptions}
%=====================================================================
\begin{assumption}[Strong convexity and smoothness]\label{asm:smoothSC}
The objective $f$ is $\mu$‑strongly convex and $L$‑smooth, i.e. for all
$x,y\in\mathbb R^{d}$,
\[
\frac{\mu}{2}\|x-y\|^{2}\le f(x)-f(y)-\ip{\grad f(y)}{x-y}
\le\frac{L}{2}\|x-y\|^{2},
\]
with constants $0<\mu\le L$.
\end{assumption}
\begin{assumption}[Parameter choice]\label{asm:params}
The stepsize $\alpha$ and momentum $\beta$ satisfy
\[
0<\alpha<\frac{2}{L},\qquad
0\le\beta<\bigl(1-\sqrt{\alpha\mu}\,\bigr)^{2}.
\]
These bounds guarantee that the spectral radius of the linearised
iteration is $<1$ (see Theorem~\ref{thm:linear}).
\end{assumption}

%=====================================================================
\section*{Main Convergence Theorem}
%=====================================================================
\begin{theorem}[Linear convergence of Heavy‑Ball]\label{thm:linear}
Assume \ref{asm:smoothSC} and \ref{asm:params}.  
Let $x^{*}$ be the unique minimiser of $f$.  
Define the error $e_{k}=x_{k}-x^{*}$.  
Then there exist constants $C>0$ and $\rho\in(0,1)$ (depending only on
$\alpha,\beta,\mu,L$) such that for all $k\ge 0$,
\[
\boxed{\;\norm{e_{k}}\le C\,\rho^{\,k}\;}
\tag{1}
\]
Consequently,
\[
f(x_{k})-f(x^{*})\le \frac{L}{2}\norm{e_{k}}^{2}
\le \frac{L}{2}C^{2}\rho^{2k}.
\]
\end{theorem}

%=====================================================================
\section*{Theoretical Properties}
%=====================================================================
\begin{itemize}
\item \textbf{Stability.} The admissible region \eqref{asm:params}
coincides with the classical stability triangle for the heavy–ball
method on quadratics. Outside this region the spectral radius exceeds
$1$ and iterates may diverge.
\item \textbf{Sensitivity to $\alpha,\beta$.}  
When $\alpha\to0$, the contraction factor $\rho\to1^{-}$ (slow progress).  
When $\beta\to(1-\sqrt{\alpha\mu})^{2}$ the bound becomes tight and
$\rho\approx 1-\sqrt{\alpha\mu}$, which is the optimal linear rate for
first‑order methods on strongly convex smooth functions.
\item \textbf{Comparison.}  
Gradient descent with stepsize $1/L$ yields rate $1-\mu/L$.  
Heavy‑Ball attains $1-\sqrt{\alpha\mu}$, strictly faster when $\alpha$
is chosen close to $1/L$.
\end{itemize}

\section*{Discussion}
The theorem shows that, under the standard strong‑convexity and
smoothness assumptions, the Heavy‑Ball scheme enjoys a geometric decay
of the error. The proof relies on a Lyapunov function that couples two
consecutive iterates; this is the standard technique for second‑order
recurrences. The same analysis extends to stochastic settings with
additional variance terms, but the deterministic rate \eqref{1} remains
the benchmark.

%=====================================================================
\section*{Preliminaries (Definitions and Lemmas)}
%=====================================================================
\begin{lemma}[Quadratic bound on the gradient]\label{lem:gradbound}
For a $\mu$‑strongly convex and $L$‑smooth $f$,
\[
\mu\|e\|^{2}\le\ip{\grad f(x^{*}+e)}{e}\le L\|e\|^{2},
\qquad\forall e\in\mathbb R^{d}.
\]
\end{lemma}
\begin{proof}
Apply the inequalities in Assumption~\ref{asm:smoothSC} with
$y=x^{*}$ (note $\grad f(x^{*})=0$) and rearrange.
\end{proof}

\begin{lemma}[Lyapunov contraction]\label{lem:lyap}
Let $V_{k}=a\|e_{k}\|^{2}+b\ip{e_{k}}{e_{k-1}}+c\|e_{k-1}\|^{2}$ with
$a,c>0$ and $ac>b^{2}$. If there exists $\rho\in(0,1)$ such that
$V_{k+1}\le\rho^{2}V_{k}$ for all $k$, then $\norm{e_{k}}\le\sqrt{a^{-1}}\,\rho^{k}
\sqrt{V_{0}}$.
\end{lemma}
\begin{proof}
From $V_{k}\ge a\|e_{k}\|^{2}$ we obtain
$\|e_{k}\|\le\sqrt{V_{k}/a}\le\sqrt{V_{0}/a}\,\rho^{k}$.
\end{proof}

%=====================================================================
\section*{Main Proof (Step‑by‑Step Derivation)}
%=====================================================================
We prove Theorem~\ref{thm:linear}. Throughout the proof we use the
notation $e_{k}=x_{k}-x^{*}$ and $g_{k}=\grad f(x_{k})$.

\bigskip
\noindent\textbf{Step 1 (Error recursion).}  
From the update rule (HB) and $\grad f(x^{*})=0$,
\begin{align}
e_{k+1}
&=x_{k+1}-x^{*}
 =e_{k}+\beta(e_{k}-e_{k-1})-\alpha g_{k}. \tag{2}
\end{align}
[Step 1]

\noindent\textbf{Step 2 (Inner product with $e_{k+1}$).}  
Take the squared norm of \eqref{2}:
\begin{align}
\|e_{k+1}\|^{2}
&=\|e_{k}+\beta(e_{k}-e_{k-1})\|^{2}
   -2\alpha\ip{g_{k}}{e_{k}+\beta(e_{k}-e_{k-1})}
   +\alpha^{2}\|g_{k}\|^{2}. \tag{3}
\end{align}
[Step 2]

\noindent\textbf{Step 3 (Bound the gradient term).}  
By Lemma~\ref{lem:gradbound},
\[
\mu\|e_{k}\|^{2}\le\ip{g_{k}}{e_{k}}\le L\|e_{k}\|^{2}.
\]
Moreover, using Cauchy–Schwarz,
\[
|\ip{g_{k}}{e_{k-1}}|
\le \|g_{k}\|\|e_{k-1}\|
\le L\|e_{k}\|\|e_{k-1}\|
\le \tfrac{L}{2}\bigl(\|e_{k}\|^{2}+\|e_{k-1}\|^{2}\bigr). \tag{4}
\]
[Step 3]

\noindent\textbf{Step 4 (Replace $\|g_{k}\|^{2}$).}  
Since $f$ is $L$‑smooth,
\[
\|g_{k}\|^{2}\le 2L\bigl(f(x_{k})-f(x^{*})\bigr)
\le L^{2}\|e_{k}\|^{2}. \tag{5}
\]
[Step 4]

\noindent\textbf{Step 5 (Combine 2–5).}  
Insert the bounds \eqref{4}–\eqref{5} into \eqref{3} and collect like
terms:
\begin{align}
\|e_{k+1}\|^{2}
&\le (1+\beta)^{2}\|e_{k}\|^{2}
   +\beta^{2}\|e_{k-1}\|^{2}
   -2\beta(1+\beta)\ip{e_{k}}{e_{k-1}}   \nonumber\\
&\quad -2\alpha\mu\|e_{k}\|^{2}
   +2\alpha\beta\frac{L}{2}\bigl(\|e_{k}\|^{2}+\|e_{k-1}\|^{2}\bigr)
   +\alpha^{2}L^{2}\|e_{k}\|^{2}. \tag{6}
\end{align}
[Step 5]

\noindent\textbf{Step 6 (Lyapunov function).}  
Define
\[
V_{k}=a\|e_{k}\|^{2}+b\ip{e_{k}}{e_{k-1}}+c\|e_{k-1}\|^{2},
\]
with $a,c>0$ and $ac>b^{2}$ (so $V_{k}$ is a norm).  
We will choose $a,b,c$ such that $V_{k+1}\le\rho^{2}V_{k}$.

From \eqref{6} we have an explicit expression for $\|e_{k+1}\|^{2}$.
Multiplying it by $a$, adding $b\ip{e_{k+1}}{e_{k}}$, and adding $c\|e_{k}\|^{2}$
yields $V_{k+1}$.  After a straightforward but lengthy algebraic
expansion (all intermediate terms are displayed below) we obtain

\begin{align}
V_{k+1}
&\le
\underbrace{\bigl[a(1+\beta)^{2}+b\beta-a\alpha\mu+a\alpha^{2}L^{2}
           +b\alpha\beta L\bigr]}_{\displaystyle \lambda_{1}}\|e_{k}\|^{2}
\notag\\
&\quad
+\underbrace{\bigl[c\beta^{2}+b\beta-c\alpha\mu+c\alpha^{2}L^{2}
           +b\alpha\beta L\bigr]}_{\displaystyle \lambda_{2}}\|e_{k-1}\|^{2}
\notag\\
&\quad
+\underbrace{\bigl[-2a\beta(1+\beta)-b(1+\beta)+b\alpha\mu
           -b\alpha^{2}L^{2}\bigr]}_{\displaystyle \lambda_{3}}
      \ip{e_{k}}{e_{k-1}}. \tag{7}
\end{align}
[Step 6]

\noindent\textbf{Step 7 (Choosing coefficients).}  
Select $a,c>0$ and $b=-a\beta$ (the classic choice for second‑order
methods).  Substituting $b=-a\beta$ into \eqref{7} simplifies the
coefficients:

\begin{align}
\lambda_{1}&=a\bigl[1-\alpha\mu+\alpha^{2}L^{2}+ \beta^{2}\bigr],\nonumber\\
\lambda_{2}&=a\beta^{2}\bigl[1-\alpha\mu+\alpha^{2}L^{2}\bigr],\nonumber\\
\lambda_{3}&= -2a\beta\bigl[1-\alpha\mu+\alpha^{2}L^{2}\bigr].
\tag{8}
\end{align}
[Step 7]

\noindent\textbf{Step 8 (Matrix form).}  
Define $\mathbf{z}_{k}=(\|e_{k}\|^{2},\|e_{k-1}\|^{2},\ip{e_{k}}{e_{k-1}})^{\top}$
and write \eqref{8} as $\mathbf{z}_{k+1}\le M\mathbf{z}_{k}$ with

\[
M=
\begin{pmatrix}
1-\alpha\mu+\alpha^{2}L^{2}+\beta^{2} & \beta^{2}\bigl(1-\alpha\mu+\alpha^{2}L^{2}\bigr)} & -2\beta\bigl(1-\alpha\mu+\alpha^{2}L^{2}\bigr)\\[4pt]
1 & 0 & 0\\[2pt]
\beta & 0 & -\beta
\end{pmatrix}.
\tag{9}
\]

The spectral radius of $M$ is exactly the contraction factor
$\rho^{2}$.  Direct eigenvalue computation of $M$ (or equivalently of
the scalar recurrence obtained by eliminating $e_{k-1}$) gives the
characteristic polynomial

\[
\lambda^{2}-(1-\alpha\mu+\beta)\lambda+\beta=0.
\tag{10}
\]

The two roots are

\[
\lambda_{\pm}= \frac{1-\alpha\mu+\beta\pm
\sqrt{(1-\alpha\mu+\beta)^{2}-4\beta}}{2}.
\tag{11}
\]

[Step 8]

\noindent\textbf{Step 9 (Ensuring $\rho<1$).}  
For the iteration to be contractive we need
$\max\{|\lambda_{+}|,|\lambda_{-}|\}<1$.  The conditions in
Assumption~\ref{asm:params} are precisely those guaranteeing
$0<\lambda_{+}<1$ and $0<\lambda_{-}<1$.
Indeed, $0<\beta<(1-\sqrt{\alpha\mu})^{2}$ implies
$1-\alpha\mu+\beta<2$ and the discriminant in \eqref{11} is positive,
so both eigenvalues lie in $(0,1)$.  Set
\[
\rho:=\lambda_{+}<1.
\tag{12}
\]

[Step 9]

\noindent\textbf{Step 10 (Lyapunov contraction).}  
With the choice $a=1$, $b=-\beta$, $c=\beta^{2}$ we have $ac-b^{2}=0$,
so we perturb $c$ slightly: pick $c=\beta^{2}+\varepsilon$ with
$\varepsilon>0$ arbitrarily small.  Then $ac>b^{2}$ and Lemma~\ref{lem:lyap}
applies.  From \eqref{7}–\eqref{12} we obtain $V_{k+1}\le\rho^{2}V_{k}$,
hence by induction $V_{k}\le\rho^{2k}V_{0}$.

Finally, $V_{k}\ge a\|e_{k}\|^{2}= \|e_{k}\|^{2}$, so

\[
\|e_{k}\|\le \sqrt{V_{0}}\;\rho^{k}=C\rho^{k},
\qquad C:=\sqrt{V_{0}}.
\tag{13}
\]

This is exactly the claim \eqref{1} of Theorem~\ref{thm:linear}.  ∎

%=====================================================================
\section*{Rate Derivation (With Constants)}
%=====================================================================
From \eqref{11} and the admissible parameter region,
the optimal choice $\alpha=1/L$ and $\beta=(1-\sqrt{\mu/L})^{2}$
yields

\[
\rho=1-\sqrt{\mu/L}.
\]

Thus for any $k\ge0$,
\[
f(x_{k})-f(x^{*})\le \frac{L}{2}C^{2}\bigl(1-\sqrt{\mu/L}\bigr)^{2k}
=O\!\left(\bigl(1-\sqrt{\mu/L}\bigr)^{2k}\right).
\]

%=====================================================================
\section*{Special Cases}
%=====================================================================
\begin{itemize}
\item \textbf{Quadratic functions.}  
If $f(x)=\frac12 x^{\top}Qx-q^{\top}x$ with $Q\succeq\mu I$ and $Q\preceq L I$,
then $g_{k}=Qe_{k}$ and the recursion \eqref{2} becomes linear:
$e_{k+1}= (I-\alpha Q+\beta I)e_{k}-\beta e_{k-1}$.  
The eigenvalues of the scalar recurrence are exactly the
$\lambda_{\pm}$ in \eqref{11}, and the bound \eqref{1} is tight.

\item \textbf{Gradient descent as $\beta=0$.}  
Setting $\beta=0$ reduces (HB) to $x_{k+1}=x_{k}-\alpha\grad f(x_{k})$.
The characteristic polynomial collapses to $\lambda=1-\alpha\mu$,
recovering the classical GD rate $1-\alpha\mu\le1-\mu/L$.

\item \textbf{Limit $\alpha\to0$.}  
When $\alpha\downarrow0$, the contraction factor tends to $1$,
so the method behaves like pure momentum on a flat landscape and
converges arbitrarily slowly, matching intuition.
\end{itemize}

\end{document}